\sloppy

%
% For taking out stuff
%
\newcommand{\takeout}[1]{\empty}

% Original draft overfull box marker omitted; warnings remain.

% Original UTF-8 source is read directly by XeLaTeX.
% XeLaTeX TU text encoding retained for original Unicode bibliography characters.
%\usepackage{microtype}

\usepackage{graphicx}
\usepackage{float}       %table placement



%\usepackage{amssymb}

\synctex=1


%\bibliographystyle{ACM-Reference-Format}

%%%%%%%%%%%%%%%%%%%%%%%%%%%%%%%%%%%%%
%%%%%%%%%%%%%%%%%%%%%%%%%%%%%%%%%%%%%
%%%%%%%%%%%%%%%%%%%%%%%%%%%%%%%%%%%%%
%%%%%%%%%%%%%%%%%%%%%%%%%%%%%%%%%%%%%Packages
%%%%%%%%%%%%%%%%%%%%%%%%%%%%%%%%%%%%%
%%%%%%%%%%%%%%%%%%%%%%%%%%%%%%%%%%%%%
%%%%%%%%%%%%%%%%%%%%%%%%%%%%%%%%%%%%%
%%%%%%%%%%%%%%%%%%%%%%%%%%%%%%%%%%%%%
\usepackage[capitalise,nameinlink]{cleveref}
%\usepackage{amssymb}
%\usepackage{MnSymbol}
%\usepackage{stmaryrd}
%\usepackage{bussproofs}
%\usepackage[usenames]{xcolor}
%\usepackage{amsfonts}
% Original UTF-8 source is read directly by XeLaTeX.
\usepackage{latexsym}
%\usepackage[noadjust]{cite}
%\usepackage{multicol}
\usepackage{mathrsfs}
%\usepackage{bbm}
\usepackage{dsfont}

\usepackage{tikz-cd}
\usetikzlibrary{arrows,arrows.meta,calc}
%
% Somehow the standard arrow head look very big and ugly here;
% so we're changing their appearance.
%
%\tikzcdset{arrow style=tikz, diagrams={>={Computer Modern Rightarrow[width=4pt,length=3pt]}}}
\tikzcdset{arrow style=tikz, diagrams={>=to}}

\tikzstyle{shiftarr}=[
        rounded corners,%
        to path={--([#1]\tikztostart.center)
                     -- ([#1]\tikztotarget.center) \tikztonodes
                     -- (\tikztotarget)},
]

\newcommand{\pullbackangle}[2][]{\arrow[phantom,to path={
                     -- ($ (\tikztostart)!1cm!#2:([xshift=8cm]\tikztostart) $)
                        node[anchor=west,pos=0.0,rotate=#2,
                        inner xsep = 0]
                        {\begin{tikzpicture}[minimum
                        height=1mm,baseline=0,#1]
    \draw[-] (0,0) -- (.5em,.5em) -- (0,1em);
                        \end{tikzpicture}}}]{}}

\usepackage{framed}
\usepackage{multicol}
\usepackage{mathtools}
\usepackage{soul}


%
% Theorem environments; needs to be loaded before hyperref so that
% \autoref works

%\newcommand{\ownthmSpaceAbove}{5pt}
%\newcommand{\ownthmSpaceBelow}{5pt}
% all thmtools hooks are within the theorem group, so we need to adjust global definitions
\newcommand{\resetCurThmBraces}{%
\gdef\curThmBraceOpen{(}%
\gdef\curThmBraceClose{)}}
\resetCurThmBraces
\newcommand{\removeThmBraces}{%
\gdef\curThmBraceOpen{}%
\gdef\curThmBraceClose{}}

\newenvironment{notheorembrackets}{\removeThmBraces}{\resetCurThmBraces}
\usepackage{etoolbox}
\patchcmd{\thmhead}{(#3)}{\curThmBraceOpen #3\curThmBraceClose }{}{}

\takeout{
\declaretheoremstyle[
    spaceabove=\ownthmSpaceAbove,
    spacebelow=\ownthmSpaceBelow,
    headpunct=.,
    postheadspace=.5em,
    notebraces={\curThmBraceOpen}{\curThmBraceClose},
    postheadhook={\resetCurThmBraces},
]{definition}
\declaretheoremstyle[
    style=definition,
    bodyfont=\itshape,
    notebraces={\curThmBraceOpen}{\curThmBraceClose},
    postheadhook={\resetCurThmBraces},
]{theorem}}

%
%Equations numbered within section
%
\numberwithin{equation}{section}



%%%%%%%%%%%%%%%%%%%%%%%%%%%%%%%%%%%%%
%%%%%%%%%%%%%%%%%%%%%%%%%%%%%%%%%%%%%
%%%%%%%%%%%%%%%%%%%%%%%%%%%%%%%%%%%%%
%%%%%%%%%%%%%%%%%%%%%%%%%%%%%%%%%%%%%HelpMe
%%%%%%%%%%%%%%%%%%%%%%%%%%%%%%%%%%%%%
%%%%%%%%%%%%%%%%%%%%%%%%%%%%%%%%%%%%%
%%%%%%%%%%%%%%%%%%%%%%%%%%%%%%%%%%%%%
%%%%%%%%%%%%%%%%%%%%%%%%%%%%%%%%%%%%%
%
%Margin notes
%
 \usepackage[author=anonymous,nomargin,marginclue,footnote,final]{fixme}
 \FXRegisterAuthor{ls}{als}{LS}
 \FXRegisterAuthor{cf}{acf}{CF}
 \FXRegisterAuthor{sm}{asm}{SM}
 \FXRegisterAuthor{bk}{abk}{BK}
 \FXRegisterAuthor{hb}{ahb}{HB}
\usepackage{siunitx}
\usepackage{hyperref}
\hypersetup{hidelinks,final,hypertexnames=false}

\usepackage{amsmath}
%
% Theorem environments; needs to be loaded before hyperref so that
% \autoref works; also autoref doesn't work without thmtools
%
\usepackage{amsthm}   %thmtools

% show sections and subsections in toc (needed for hyperref)
%\setcounter{tocdepth}{2}

%\addto\extrasbritish{% only necessary if babel is used
\renewcommand{\sectionautorefname}{Section}%
\renewcommand{\subsectionautorefname}{Section}%
\renewcommand{\subsubsectionautorefname}{Section}%
  %\renewcommand{\sectionname}{Section}%
%}

%
% Custom enumerate environment
%
\usepackage[inline]{enumitem}
\setlist[enumerate,1]{label=(\arabic*),font=\normalfont,align=left,leftmargin=0pt,labelindent=0pt,
listparindent=\parindent,labelwidth=0pt,itemindent=!,topsep=3pt,parsep=0pt,itemsep=3pt,start=1}
\setlist[enumerate,2]{label=(\alph*),font=\normalfont,labelindent=*,leftmargin=*,topsep=3pt,start=1}
\setlist[itemize]{labelindent=*,leftmargin=*,topsep=5pt,itemsep=3pt}
\setlist[description]{labelindent=*,leftmargin=*,itemindent=-1 em}

%
% Nice display of labels
%
%\usepackage[inline]{showlabels}
%\showlabels{bibitem}

% \usepackage{seqsplit}
% \usepackage{xstring}
% \usepackage{xcolor}
% \usepackage{lscape}

% \newcommand{\blue}[1]{\protect{\color{blue}{#1}}}
% \def\red{\color{red}}
% \def\blk{\color{black}}

% \makeatletter
% \renewcommand*{\showlabelsetlabel}[1]{%
%   %\noexpandarg%
%   \StrSubstitute{#1}{_}{\textunderscore}[\TEMP]%
%   %\parbox[t]{1.25cm}{\raggedright\normalfont\small\color{red!50!black}\showlabelfont
%   %\expandafter\seqsplit\expandafter{\TEMP}} % does not look so good
%   {\normalfont\small\color{red!50!black}\showlabelfont\{\TEMP\}}
% }
% \makeatother

\takeout{%% Alternative to nice display of labels which does not looks
         %% so good with two-column style
%
% Showkeys
%
%\usepackage[notref,notcite]{showkeys}

%%%% Text wrapping for showkeys
\usepackage{seqsplit}
\usepackage{xstring}
\newcommand{\defaultshowkeysformat}[1]{%
% instead of \textvisiblespace you can also put in ~
% if you want to keep a plain space at space characters
\StrSubstitute{#1}{ }{\textvisiblespace}[\TEMP]%
\parbox[t]{\marginparwidth}{\raggedright\normalfont\small\ttfamily\(\{\){\color{red!50!black}\expandafter\seqsplit\expandafter{\TEMP}}\(\}\)}%
}
\newenvironment{hideshowkeys}{%
  \renewcommand*\showkeyslabelformat[1]{%
  \noexpandarg%
  }
}{%
  \renewcommand*\showkeyslabelformat[1]{%
  \noexpandarg%
  \defaultshowkeysformat{##1}%
  }
}

\renewcommand*\showkeyslabelformat[1]{%
\noexpandarg%
\defaultshowkeysformat{#1}%
}
}% end takeout

%
%Equations numbered within section
%
\numberwithin{equation}{section}



%%%%%%%%%%%%%%%%%%%%%%%%%%%%%%%%%%%%%
%%%%%%%%%%%%%%%%%%%%%%%%%%%%%%%%%%%%%
%%%%%%%%%%%%%%%%%%%%%%%%%%%%%%%%%%%%%
%%%%%%%%%%%%%%%%%%%%%%%%%%%%%%%%%%%%%Environments
%%%%%%%%%%%%%%%%%%%%%%%%%%%%%%%%%%%%%
%%%%%%%%%%%%%%%%%%%%%%%%%%%%%%%%%%%%%
%%%%%%%%%%%%%%%%%%%%%%%%%%%%%%%%%%%%%
%%%%%%%%%%%%%%%%%%%%%%%%%%%%%%%%%%%%%

\theoremstyle{plain}
\newtheorem{thm}{Theorem}[section]
\newaliascnt{lem}{thm}
\newtheorem{lem}[lem]{Lemma}
\aliascntresetthe{lem}
\crefname{lem}{Lemma}{Lemmas}
\newaliascnt{propn}{thm}
\newtheorem{propn}[propn]{Proposition}
\aliascntresetthe{propn}
\crefname{propn}{Proposition}{Propositions}
\newaliascnt{cor}{thm}
\newtheorem{cor}[cor]{Corollary}
\aliascntresetthe{cor}
\crefname{cor}{Corollary}{Corollarys}
\newtheorem{target}{Target Theorem}



\theoremstyle{definition}
\newaliascnt{defn}{thm}
\newtheorem{defn}[defn]{Definition}
\aliascntresetthe{defn}
\crefname{defn}{Definition}{Definitions}
\newaliascnt{expl}{thm}
\newtheorem{expl}[expl]{Example}
\aliascntresetthe{expl}
\crefname{expl}{Example}{Examples}
\newaliascnt{constrn}{thm}
\newtheorem{constrn}[constrn]{Construction}
\aliascntresetthe{constrn}
\crefname{constrn}{Construction}{Constructions}
\newaliascnt{rem}{thm}
\newtheorem{rem}[rem]{Remark}
\aliascntresetthe{rem}
\crefname{rem}{Remark}{Remarks}
\newaliascnt{notn}{thm}
\newtheorem{notn}[notn]{Notation}
\aliascntresetthe{notn}
\crefname{notn}{Notation}{Notations}
\newaliascnt{conv}{thm}
\newtheorem{conv}[conv]{Convention}
\aliascntresetthe{conv}
\crefname{conv}{Convention}{Conventions}
\newaliascnt{conj}{thm}
\newtheorem{conj}[conj]{Conjecture}
\aliascntresetthe{conj}
\crefname{conj}{Conjecture}{Conjectures}
%\newcounter{assumption}
\newaliascnt{assn}{thm}
\newtheorem{assn}[assn]{Assumption}
\aliascntresetthe{assn}
\crefname{assn}{Assumption}{Assumptions}






%%%%%%%%%%%%%%%%%%%%%%%%%%%%%%%%%%%%%
%%%%%%%%%%%%%%%%%%%%%%%%%%%%%%%%%%%%%
%%%%%%%%%%%%%%%%%%%%%%%%%%%%%%%%%%%%%
%%%%%%%%%%%%%%%%%%%%%%%%%%%%%%%%%%%%%Commands
%%%%%%%%%%%%%%%%%%%%%%%%%%%%%%%%%%%%%
%%%%%%%%%%%%%%%%%%%%%%%%%%%%%%%%%%%%%
%%%%%%%%%%%%%%%%%%%%%%%%%%%%%%%%%%%%%
%%%%%%%%%%%%%%%%%%%%%%%%%%%%%%%%%%%%%

% From Stefan

%
% The dual of colon is lon, and the dual of the coefficient ring is ...
%
\newcommand*\lon{%
        \nobreak
        \mskip6mu plus1mu
        \mathpunct{}%
        \nonscript
        \mkern-\thinmuskip
        {:}%
        \mskip2mu
        \relax
}%

%\newcommand{\ol}[1]{\overline{#1}}
\newcommand{\overbar}[1]{\mkern 1.5mu\overline{\mkern-2.5mu#1\mkern-1.5mu}\mkern 1.5mu}
%
% Usage \mybar{<scale factor>}{<right shift>}{<math expression>}
% This macro doesn't work in tikz-cd labels
%
%\newcommand{\mybar}[3]{%
%  \mathrlap{\hspace{#2}\overline{\scalebox{#1}[1]{\phantom{\ensuremath{#3}}}}}\ensuremath{#3}
%}
\newcommand{\mybar}[3]{%
  \mathrlap{\hspace{#2}\overline{\scalebox{#1}[1]{\hphantom{\ensuremath{#3}}\vphantom{\rule{0pt}{5.5pt}}}}}\ensuremath{#3}
}

\DeclareMathOperator{\Ker}{\mathsf{ker}}

\newcommand{\infgameeq}{\sim_{\mathcal{G}}}
\newcommand{\barM}{\mybar{0.8}{2pt}{M}_1}
\newcommand{\barN}{\mybar{0.8}{2pt}{N}_1}
\newcommand{\barZ}{\mybar{0.7}{1.75pt}{Z}}

\DeclareMathOperator{\Fun}{\mathsf{Fun}}

\newcommand{\eps}{\varepsilon}
\newcommand{\syn}{\mathsf{syn}}

\newcommand{\id}{\mathsf{id}}
\newcommand{\subto}{\hookrightarrow}
\newcommand{\epito}{\twoheadrightarrow}
\newcommand{\monoto}{\rightarrowtail}
\newcommand{\inl}{\mathsf{inl}}
\newcommand{\inr}{\mathsf{inr}}
\newcommand{\inj}{\mathsf{in}}
\newcommand{\xra}[1]{\xrightarrow{~#1~}}
\newcommand{\xla}[1]{\xleftarrow{~#1~}}

\newcommand{\ola}{\overleftarrow}

\newcommand{\fair}[1]{\langle #1\rangle}




%monad/theory matters
\newcommand{\M}{\mathbb{M}}                   %monad M
\newcommand{\T}{\mathcal{T}}                     %monad/theory T
\newcommand{\arity}{\mathsf{ar}}





%categories
\newcommand{\CatC}{\mathscr{C}}           %general category C/D
\newcommand{\CatD}{\mathscr{D}}
\newcommand{\CatK}{\mathscr{K}}
\newcommand{\CatI}{\mathscr{I}}
\newcommand{\CatG}{\mathscr{G}}
\newcommand{\K}{\mathscr{K}}
\newcommand{\Met}{\mathbf{Met}}
\newcommand{\Set}{\mathbf{Set}}
\newcommand{\Pos}{\mathbf{Pos}}
\newcommand{\Mon}{\mathscr{M}}             %general monoidal category
\newcommand{\bfC}{\mathbf{C}}                    %explictly defined category
\newcommand{\bfD}{\mathbf{D}}                    %domain of a diagram
\DeclareMathOperator{\Alg}{\mathsf{Alg}}



% relations/operations
\setul{-1pt}{.4pt}
%\newcommand{\bisim}{\mathbin{\mathrlap{\hspace{1pt}\rule[-.5pt]{8pt}{.4pt}}\mathord{\leftrightarrow}}}
\DeclareMathOperator{\bisim}{\mathrlap{\hspace{1pt}\rule[-.3pt]{8pt}{.5pt}}\mathord{\leftrightarrow}}
\DeclareMathOperator{\similar}{\mathrlap{\hspace{.5pt}\rule[-.3pt]{8.5pt}{.5pt}}\mathord{\rightarrow}}
%\DeclareMathOperator{\trace}{\equiv_{\mathcal{T}}}
\DeclareMathOperator{\Coalg}{\mathsf{Coalg}}
\DeclareMathOperator{\beh}{\mathsf{beh}}
%\DeclareMathOperator{\Trace}{\mathcal{T}}
\DeclareMathOperator{\Support}{\mathsf{supp}}
\DeclareMathOperator{\Terms}{\mathscr{T}}
\newcommand{\Termsarg}[1]{\Terms_{\!#1}}
%\DeclareMathOperator{\TermsSi}{\Terms_{\!\Sigma,\infty}}
%\DeclareMathOperator{\TermsS1}{\Terms_{\!\Sigma,1}}
\DeclareMathOperator{\Win}{\mathsf{Win}}

\newcommand{\sem}[1]{\llbracket #1 \rrbracket}% Interpretation of modalities in S-game

%number sets/objects
\newcommand{\N}{\mathds{N}}            %natural numbers

\newcommand{\A}{\mathcal{A}}            %action set
\newcommand{\B}{\mathcal{B}}
\newcommand{\CalG}{\mathcal{G}}             %game G
\newcommand{\bfx}{\textbf{x}}                %vector
\DeclareMathOperator{\colim}{\mathsf{colim}}    %colim

\newcommand{\F}{\mathscr{F}}
\newcommand{\CalS}{\mathcal{S}}        %graded semantics abbreviation
\newcommand{\bbD}{\mathbb{D}} %Duplicator
\newcommand{\bbS}{\mathbb{S}}       %Spoiler

\newcommand{\E}{\mathcal{E}}   %equations

\DeclareMathOperator{\JSL}{\mathsf{JSL}}

\newcommand{\CalGp}{\mathcal {G}'}
\newcommand{\PT}{\mathsf{PT}}
%\newcommand{\Coef}{\text{Coef}}
\newcommand{\Smds}{\mathsf{Smds}}
\newcommand{\dl}{\mathbf 0}


%functors
\newcommand{\Id}{\mathsf{Id}}                 %identity
\newcommand{\pow}{\mathscr P}              %powerset
\newcommand{\powf}{\pow_{\mathsf f}}
\newcommand{\CalN}{\mathcal{N}}           %neighborhood
\newcommand{\CalD}{\mathcal{D}}            %distribution
\newcommand{\hra}{\hookrightarrow}         %inclusion
\newcommand{\cpow}{\ola{\pow}}               %contravariant powerset functor
\newcommand{\Q}{\mathcal{Q}}                  %contra powerset functor
\newcommand{\bbC}{\mathcal{C}}              %canonical M_1-algebra

\usepackage{booktabs}   %% For formal tables:
                        %% http://ctan.org/pkg/booktabs
\usepackage{subcaption} %% For complex figures with subfigures/subcaptions
                        %% http://ctan.org/pkg/subcaption
